%! TeX root: ../informe.tex
\section{Estado de la cuestión}
\label{sec:art}

\subsection{Análisis automático de grafiti}
\label{subsec:art-grafiti}

La literatura sobre la aplicación de visión por computador y agrupamiento no
supervisado al análisis de grafiti es escasa, reciente y de alcance acotado.
\citea[etali]{Fogaça}{fogaca_deep_2023} presentan un sistema encargado por el
Ayuntamiento de Lisboa para clasificar imágenes de calle en grafiti ilegal, arte
urbano o ausencia de grafiti, y localizarlo mediante un detector
\textit{fine-tuned}. \citea[etali]{Tokuda}{tokuda_quantifying_2019} cuantifican
la densidad de grafitis en São Paulo a partir de Google Street View con un
modelo Mask-RCNN entrenado sobre anotaciones propias. Ambos establecen la
viabilidad del aprendizaje profundo sobre imágenes captadas en condiciones no
controladas. En el ámbito del agrupamiento, aunque varios trabajos categorizan
grafiti por tipo o estilo, la mayoría lo hacen de forma
manual~\cite{taylor_twenty-first_2010};
\citea[etali]{Garcia~Garcia}{garcia_garcia_graffiti_2025} proponen un
agrupamiento basado en los colores dominantes de cada pieza (K-medias sobre RGB
y Lab), pero prescinden de la estructura espacial, la textura y el contenido
semántico. La tabla~\ref{tab:art-comparativa} sitúa estos trabajos frente al
presente: ninguno combina descriptores visuales profundos con técnicas estándar
de agrupamiento ni acompaña la evaluación de una caracterización del coste, que
es la carencia que este trabajo pretende cubrir.

{\scriptsize
\setlength{\tabcolsep}{3pt}
\begin{longtable}{|p{1.9cm}|p{2.1cm}|p{2.0cm}|p{2.2cm}|p{2.2cm}|p{1.9cm}|} \hline
  \textbf{Trabajo} & \textbf{Tarea} & \textbf{Datos} & \textbf{Características} & \textbf{Técnica} & \textbf{Métrica} \\ \hline
  \endhead
  Fogaça (2023) & Clasificación y detección & Imágenes de calle (Lisboa) & CNN \textit{fine-tuned} & Clasificación + detector & Exactitud / mAP \\ \hline
  Tokuda (2019) & Cuantificación de densidad & Street View (São Paulo) & Mask-RCNN & Segmentación supervisada & Precisión / \textit{recall} \\ \hline
  Taylor (2010) & Categorización por contenido & Fotografías de grafiti & Inspección visual & Categorización manual & --- \\ \hline
  García~García (2025) & Agrupamiento de grafitis & Imágenes de grafiti & Colores dominantes (RGB / Lab) & K-medias & Interna \\ \hline
  \textbf{Este trabajo} & Agrupamiento y recuperación & Banco de grafiti & \textit{Embeddings} profundos & Reducción + \textit{clustering} clásico & Interna, externa y coste \\ \hline
  \caption[Comparativa de trabajos previos sobre análisis de grafiti]{Comparativa de trabajos previos frente al planteamiento de este trabajo.}
  \label{tab:art-comparativa}
\end{longtable}
}

\subsection{Similitud, agrupamiento e identificación de autoría}
\label{subsec:art-emb-red-clus}

El agrupamiento de imágenes encadena habitualmente tres etapas: extracción de
una representación vectorial, reducción de dimensionalidad y agrupamiento. Los
modelos de clasificación preentrenados como ResNet~\cite{he_deep_2015} se usan
ampliamente como extractores, y modelos más recientes como
CLIP~\cite{radford_learning_2021} o DINOv2~\cite{oquab_dinov2_2024} ofrecen
representaciones transferibles sin ajuste de dominio. La reducción previa (PCA,
UMAP~\cite{mcinnes_umap_2020}) mitiga la inestabilidad de las distancias en alta
dimensión, y las familias clásicas de agrupamiento (centroides, densidad,
jerárquicos, probabilísticos) cubren distintos supuestos
estructurales~\cite{ester_density-based_1996,campello_density-based_2013,arthur_k-means_2007}.
Para la recuperación por similitud, los índices ANN como
HNSW~\cite{malkov_efficient_2018} o FAISS~\cite{johnson_billion-scale_2017} son
el estándar \textit{de facto}. Estos componentes están bien caracterizados de
forma individual, pero no se han localizado trabajos que los evalúen de forma
combinada sobre imágenes de grafiti ni que caractericen el coste por etapa en
función del tamaño del banco.

La identificación de autoría a partir de firmas es análoga a otros problemas ya
estudiados. Fiel y Sablatnig~\cite{fiel_writer_2015} recuperan la identidad de
un escritor a partir de las activaciones de una CNN, resolviendo la consulta por
vecindad más cercana: el patrón de \enquote{codificar una vez y recuperar por
similitud} es precisamente el que aplica este trabajo. En reconocimiento facial,
FaceNet~\cite{schroff_facenet_2015} aprende un espacio métrico mediante una
pérdida por tripletes; la autoría en grafiti puede beneficiarse de este enfoque,
al ser también un problema de aprendizaje métrico con número de identidades
desconocido \textit{a priori} y pocas muestras por clase. Una línea más reciente
integra representación y agrupamiento en un mismo entrenamiento
(\textit{deep clustering})~\cite{min_survey_2018,caron_deep_2019,van_gansbeke_scan_2020},
pero presupone un volumen de datos y de cómputo del que no se dispone aquí. Por
ello este trabajo se sitúa deliberadamente en el enfoque desacoplado:
aprovechar representaciones preentrenadas, posiblemente ajustadas con una cabeza
ligera de dominio, y combinarlas con algoritmos de \textit{clustering} maduros.
